% executive_summary.tex -- one page, standalone.
% Build:  tectonic executive_summary.tex   (run from reports/executive_summary/)
% Every figure here is quoted from the memo; nothing is computed in this file.

\documentclass[10pt,letterpaper]{article}

\usepackage[T1]{fontenc}
\usepackage[utf8]{inputenc}
\usepackage[letterpaper,margin=0.78in]{geometry}
\usepackage{newtxtext}
\usepackage{newtxmath}
\usepackage{booktabs}
\usepackage{microtype}
\usepackage[hidelinks]{hyperref}
\usepackage{titlesec}
\usepackage{enumitem}

\titleformat{\section}{\normalfont\normalsize\bfseries}{}{0em}{}
\titlespacing*{\section}{0pt}{1.3ex plus 0.3ex}{0.5ex}
\setlength{\parindent}{0pt}
\setlength{\parskip}{0.55ex plus 0.2ex}
\setlist[itemize]{leftmargin=1.1em,itemsep=0.25ex,topsep=0.3ex,parsep=0pt}
\renewcommand{\arraystretch}{1.05}
\pagestyle{empty}

\hypersetup{
  pdftitle={Executive summary: distributed sterile-injectable manufacturing under real-world constraints},
  pdfauthor={Rikhin Kavuru},
  pdfsubject={Constraint-driven feasibility investigation. Illustrative inputs.
              Not a claim of regulatory compliance or manufacturing readiness.}
}

\begin{document}

{\large\bfseries Distributed sterile-injectable manufacturing under real-world constraints}\\[0.3ex]
{\small Executive summary. Rikhin Kavuru. Full memo and code:
\texttt{reports/memo/memo.pdf}, \texttt{docs/}, \texttt{src/}.}

\vspace{0.6ex}
\fbox{\begin{minipage}{\dimexpr\textwidth-2\fboxsep-2\fboxrule\relax}
\footnotesize
\textbf{PRELIMINARY AND ILLUSTRATIVE. NOT FOR EXTERNAL CLAIMS.} 87 of 91 model parameters are
tier-5 placeholders with no operational, official, peer-reviewed or elicited source. Zero of a
target 25 to 30 interviews and zero of three required independent reviews are complete. All
sixteen regulatory gates are UNCERTAIN, so eligibility is NO\_CONCLUSION for every architecture.
Nothing here claims any system is FDA-compliant, clinically deployable, or ready to manufacture
drugs. Every number is the behaviour of a model under those inputs.
\end{minipage}}

\section{The question}
Under what product, demand, regulatory and operating conditions could regional manufacturing
reduce sterile-injectable shortages more effectively than safety stock, dual sourcing, or
additional centralized capacity? Two non-controlled small-molecule presentations were carried
through: sodium bicarbonate 8.4\% in a 50\,mL vial, in prolonged shortage, and norepinephrine
bitartrate 1\,mg/mL in a 4\,mL vial, resolved.

\section{Method, in one paragraph}
A protocol frozen before any run fixes the service target: mean fill rate $\geq 0.99$ in at least
90\% of simulated years. A daily discrete-event engine simulates twenty supply architectures over
five years on common random numbers, so every architecture meets the same disruptions in the same
years. Each is optimized to that one target, then re-evaluated at 100 paired runs. An ablation
harness over 22 named failure mechanisms attributes the result and, in practice, detects defects:
a factor whose sign comes out wrong is usually a bug. Twenty-five were found that way.

\section{What the model says}
\begin{itemize}
\item \textbf{A capacity-short network cannot be stocked out of the deficit.} On the short product,
a full year of safety stock with daily base-stock review leaves the status quo at mean fill 0.869
and a tail probability of zero. At twice base demand, none of the twenty architectures meets the
target anywhere. That region argues for capacity, not for \emph{regional} capacity.

\item \textbf{Owned regional nodes lose the years before they exist.} On the short product the
distributed design needs commissioning within 741 days, supplier disruption at or below 0.127
against an assumed base of 0.200, shelf life at or above 17.6 months and changeover within 4.5
days. It therefore sits outside its own feasible region on the supplier axis at the base
assumption, at the highest cost in the battery: 32.0 against 19.4 million USD per year for the
cheapest cell that meets the target.

\item \textbf{Adding a validated release layer to those nodes buys nothing.} It returns mean fill
identical to four decimals while costing 1.0 to 1.7 million USD per year more. The reasons are a
model defect (MD-18) and a regulatory gate held UNCERTAIN (G15), neither an argument for the layer.

\item \textbf{Regional manufacturing is not eliminated, but the honest form is narrow.} On the
product that is \emph{not} capacity-short, the distributed design meets the target at a tail
probability of 0.940, at about 9\% above rotating contracted campaigns and 32\% above the cheapest
feasible cell. It clears the bar where the incumbent plant has slack, and not where shortages
actually arise.

\item \textbf{Two architectures meet the target on both products, and neither is a plant the
sponsor builds.} Rotating prequalified campaigns across three already-registered third-party fill
lines (S17), and dual finished-dose sourcing with an explicit key-starting-material tier (S9).
\end{itemize}

\section{The conditional answer}
On these inputs, the strongest intervention is contracted registered capacity plus positioned
finished-goods stock, bought rather than built. The binding constraint is a contract, not a
machine. Of 40 re-evaluated cells, 14 meet the target and 16 sit within one standard error of it,
so two fifths of the verdicts cannot be resolved at this sample size.

\section{What would overturn this, and what it would take}
The three inputs that most change the answer are the regional replenishment rule, the rate at which
one event removes several nominally independent sites, and the real backorder window. All three are
currently placeholders the author wrote. None can be settled by more computation; each needs a
practitioner. The study also retracted three of its own headline conclusions after correcting its
own model, and one earlier failure of the distributed design proved to be an artifact of an unequal
search space rather than a property of distribution.

\end{document}
